\speciestitle{Cloud Ice Water Content}{IWC}{IWC}{%
\item[Units:] g/m\cp3
\item[Useful range:] 215\,--\,83\,hPa
\item[Contact: ] Alyn Lambert, \email{Alyn.Lambert@jpl.nasa.gov}}

\subsection{Introduction}

The MLS IWC is retrieved from cloud-induced radiances (\Tcir) of the
240-GHz window channel in a separate processing step after the
atmospheric state (Temperature and tangent height pressure) and important
gaseous species (H\cs2O, O\cs3, HNO\cs3) have been finalized in the
retrieval processing.  The derived \Tcir\ are binned onto the standard
horizontal (1.5\degsym\ along track) and vertical (12 surfaces per
decade change in pressure) grids, and converted to IWC using the modeled
\Tcir\,--\,IWC relations \citep{WuEtal:2006}.  The standard IWC profile has
a useful vertical range between 215\,--\,83\,hPa although the validation
has been conducted for a subset of the range of IWC values. IWC
measurements beyond the value ranges specified in
Table~\ref{tab:IWCSummary} are to be regarded currently as giving only
qualitative information on cloud ice.  They require further validation
for quantitative interpretation.

\subsection{Resolution}

In the IWC ranges specified in Table~\ref{tab:IWCSummary}, each MLS
measurement can be quantitatively interpreted as the average IWC for the
volume sampled. This volume has a vertical extent of \texttildelow3\,km, with
\texttildelow300\,km and 7\,km along and cross track respectively.

\subsection{Precision}

The precision values quoted in the IWC files do not represent the true
precision of the data.  The precision for a particular measurement must
be evaluated on a daily basis using the method described in the
screening section below.  The precision listed in
Table~\ref{tab:IWCSummary} reflects typical values obtained from the
method described below.

\subsection{Accuracy}

The IWC accuracy values listed in Table~\ref{tab:IWCSummary} are
estimates from comparisons of the earlier v2.2 MLS data product with
CloudSat and detailed analyses on the v2.2 error budget can be found in
\citet{WuEtal:2008}.

\subsection{Data screening}
\label{sec:IWCScreening}
\begin{description}
\item[Pressure range (215\,--\,83\,hPa):] \standardrangestatement\ The
  maximum detectable IWC is \texttildelow100\,mg/m\cp3.
\item[Use Temperature Status, Quality and Convergence:] The user is
  recommended to screen the IWC data using the \texttt{Status} field in
  the collocated temperature profile to exclude bad retrievals
  \citep{SchwartzEtal:2008:Validation}. In other words, only IWC profiles for which
  temperature \texttt{Status} is an even number should be used.
  Similarly, the users should only use IWC profiles where the
  corresponding Temperature profiles have \emph{Quality} of 0.9 or
  larger and \texttt{Convergence} less than 1.03.
\item[Other screening:] The IWC product derives from differences between
  measured radiances and those predicted assuming cloud free conditions.
  Spectroscopic and calibration uncertainties give rise to temporally
  and geographically varying biases in this difference, and hence the
  IWC product.  These biases must be iteratively identified and removed,
  using a ``$2\sigma$\,--\,$3\sigma$'' screening method, as described
  below.

\begin{enumerate}
\item Uncertainties in spectroscopy and atmospheric composition are
  manifested as residual biases in the IWC fields which should be
  identified and removed as follows.  IWC data should be averaged in a
  10\degsym\ latitude bins and outliers rejected iteratively by
  excluding measurements greater than $2\sigma$ standard deviation about
  the mean ($\mu$) of the bin. Repeat the $\sigma$ and $\mu$
  calculations after every new set of rejections.  Convergence is
  usually reached within 5\,--\,10 iterations, and the final $\sigma$ is
  the estimated precision for the IWC measurements.
\item Interpolate the final $\sigma$ and $\mu$ to the latitude of
  each measurement, and subtract $\mu$ from IWC for each measurement.
\item Finally, apply the 3$\sigma$ threshold to determine if an IWC
  measurement is statistically significant. In other words, it must have
  $\text{IWC} > \mu + 3\sigma$ in order to be considered as a
  significant cloud hit.  The $3\sigma$ threshold is needed for cloud
  detection since a small percentage of clear-sky residual noise can
  result in a large percentage of ``false alarms'' in cloud detection.
\end{enumerate}
\end{description}

\subsection{Artifacts}

At wintertime mid-to-high latitudes, strong stratospheric gravity waves
may induce large fluctuations in the retrieved tangent pressure, and
cause false cloud detection with the 2$\sigma$\,--\,3$\sigma$ screening
method. The false cloud detection seems to affect the 100\,hPa pressure
level most, as expected for such impact coming from the lower
stratosphere.

\subsection{Comparisons with other datasets}
  
Scatter plots of \thisv vs v2.2 IWC show mean biases of less than 10\%
in the pressure range 215\,--\,83\,hPa and histograms show that
the random noise in \thisv IWC is generally larger than in v2.2 (see
Figure~\ref{fig:IWC} and Table~\ref{tab:IWCSummary}). However, in
localized regions, there can be substantial differences in IWC values
between these versions that are driven by the reaction of the IWC
retrieval to the changes in \thisv UTLS H\cs2O.  These effects are largest at
pressures of 215~hPa and higher and have been verified by examining
the negative correlations between the changes in H\cs2O and IWC.

Apart from the differences noted above,
the MLS \thisv IWC is similar to the MLS v2.2 product described and
validated in \citet{WuEtal:2008}.  A revised validation paper for IWC is
not planned in the near future and users are encouraged to read
\citet{WuEtal:2008} for more information.

Comparisons between v2.2 MLS and CloudSat IWC showed good agreement with
PDF differences $<$50\% for the IWC ranges specified in
Table~\ref{tab:IWCSummary}. Comparisons with AIRS, OMI and MODIS suggest
that MLS cloud tops are slightly higher by \texttildelow1\,km than the
correlative data in general.

\begin{table}
  \caption{Summary of MLS \thisv IWC precision, accuracy, and resolution.}
\vspace{\tablecaptionsep}
\label{tab:IWCSummary}
\begin{minipage}{\linewidth} \centering\begin{tabular}{cccccc} \toprule
    \parbox{6em}{\bfseries\centering Pressure / hPa} &
    \parbox{6em}{\bfseries\centering Resolution\footnote{The along-track,
        cross-track and vertical extent, respectively of the atmospheric volume
        sampled by an individual MLS measurement.} / km} &
    \parbox{6em}{\bfseries\centering Typical precision\footnote{These are
        typical 1$\sigma$ precisions where the better values are for the
        extratropics and the poorer values for the tropics. The precision for a
        particular measurement must be evaluated on a daily basis using the
        method described in the text.} / mg/m\cp3} &
    \multicolumn{2}{c}{\parbox{12em}{\null\hfill\bfseries
        Accuracy\footnote{Estimated from comparisons with CloudSat.} /
        mg/m\cp3\hspace{8ex}\hfill\null\\%
        \null\hspace{-3ex}$<$10\,mg/m\cp3\hfill$>$10\,mg/m\cp3\hspace{4ex}\null}}
    &
    \parbox{6em}{\bfseries\centering Valid IWC range\footnote{This is the
        range where the stated precision, accuracy and resolution are
        applied. In this range MLS measurements can be quantitatively
        interpreted as the average IWC for the volume sampled. IWC values above
        this range, currently giving qualitative information on cloud ice,
        require further validation for quantitative interpretation.} \\ /
      mg/m\cp3} \\ \midrule
p$<$70 & \multicolumn{5}{c}{Unsuitable for scientific use} \\
 83 & 200$\times$7$\times$5 & 0.06\,--\,0.15     & 100\% & ---   & 0.02\,--\,50 \\
100 & 200$\times$7$\times$5 & 0.08\,--\,0.18     & 100\% & 150\% & 0.02\,--\,50 \\
121 & 250$\times$7$\times$4 & 0.08\,--\,0.24     & 100\% & 100\% & 0.04\,--\,50 \\
147 & 300$\times$7$\times$4 & 0.20\,--\,0.65 & 100\% & 100\% & 0.1\,--\,50  \\
177 & 300$\times$7$\times$4 & 0.5\,--\,2.3   & 150\% & 100\% & 0.3\,--\,50  \\
215 & 300$\times$7$\times$4 & 1.2\,--\,5.5   & 300\% & 100\% & 0.6\,--\,50  \\
p$>$260 & \multicolumn{5}{c}{Unsuitable for scientific use} \\ \bottomrule
\end{tabular}
\end{minipage}
\end{table} 
 
\begin{figure}
  \includegraphics[width=\textwidth]{figures/IWC-240-v05-00-dataqual-JUL2009-v2}
  \caption{MLS v4, v3 and v2 IWC comparisons for July 2009 at 146\,hPa
    and 100\,hPa.  (a) Left: Probability density functions (PDF) (v5,
    red; v4, blue; and v2, green) with dashed lines showing the
    corresponding noise levels (obtained by folding the negative IWC
    values about the origin) and the thin black lines representing the
    gaussian error function.  (b) Right: Scatter plots of IWC v4 vs.\ v2
    (black points) with dashed red lines indicating the 1:1 line, dashed
    yellow lines the $1\sigma$ uncertainties and the solid yellow lines are linear
    fits to the data (which lie practically on the 1:1 line).}
\label{fig:IWC}
\end{figure}

%%% Local Variables: 
%%% mode: latex
%%% TeX-master: "v5-x-report"
%%% End: 
